\documentclass[11pt]{article}
\usepackage[margin=1in]{geometry}
\usepackage{amsmath,amssymb,amsthm}
\usepackage[hidelinks]{hyperref}
\newcommand{\Tr}{\operatorname{tr}}
\title{A positive definite counterexample to the\\
nonlinear quantum-measurement monotone\\
\large Complete solution of MathDB 359456}
\author{Principia Math in Progress}
\date{12 September 2026}
\begin{document}
\maketitle

\paragraph{Task 1. Resolve the nonlinear monotone conjecture.}
A blockwise probability vector $\boldsymbol P=(P_1,\ldots,P_n)^{\mathsf T}$
consists of positive semidefinite $d\times d$ matrices summing to $I_d$.
A blockwise bistochastic matrix $\boldsymbol B=(B_{ij})$ has positive
semidefinite blocks, with every block row and block column summing to
$I_d$. The source's sequential product is
\[
 Q_i=(\boldsymbol B*\boldsymbol P)_i
       =\sum_{j=1}^n P_j^{1/2}B_{ij}P_j^{1/2}.
\]
Conjecture 1 of Rico and \.{Z}yczkowski asks whether, for every ordering
of the $Q_i$, some ordering of the $P_i$ makes
\[
 \left\|\sum_{i=1}^k(P_i-I_d/n)\right\|_2
 \geq
 \left\|\sum_{i=1}^k(Q_i-I_d/n)\right\|_2
 \quad(1\leq k\leq n),
 \qquad \|X\|_2^2=\Tr(X^*X).
\]
The answer is negative, already for $n=d=3$ and diagonal positive
definite effects and blocks.

Let
\[
 P_1=\frac1{12}\operatorname{diag}(10,1,1),\quad
 P_2=\frac1{12}\operatorname{diag}(1,10,1),\quad
 P_3=\frac1{12}\operatorname{diag}(1,1,10),
\]
and take
\[
 \boldsymbol B=
 \begin{pmatrix}
  P_1&P_2&P_3\\
  P_2&P_3&P_1\\
  P_3&P_1&P_2
 \end{pmatrix}.
\]
Each $P_j$ is positive definite, and $P_1+P_2+P_3=I_3$.
Every block row and column of $\boldsymbol B$ is a permutation of
this triple. Thus all hypotheses hold, with all nine blocks positive
definite as well.

All matrices involved are diagonal, so
$P_j^{1/2}B_{ij}P_j^{1/2}=P_jB_{ij}$. Direct multiplication yields
\[
 Q_1=P_1^2+P_2^2+P_3^2=\frac{17}{24}I_3,\qquad
 Q_2=Q_3=P_1P_2+P_2P_3+P_3P_1=\frac7{48}I_3.
\]
These output effects are also positive definite and sum to $I_3$.

Consider the displayed output ordering, with $Q_1$ first.
For every $j=1,2,3$, the diagonal entries of $P_j-I_3/3$
are a permutation of $(1/2,-1/4,-1/4)$, whence
\[
 \|P_j-I_3/3\|_2^2
 =\frac14+\frac1{16}+\frac1{16}
 =\frac38=\frac{24}{64}.
\]
In contrast,
\[
 \|Q_1-I_3/3\|_2^2
 =\left\|\frac38I_3\right\|_2^2
 =\frac{27}{64}>\frac{24}{64}.
\]
No choice of the first input effect can satisfy the required inequality
at $k=1$. Therefore no input ordering works for this output ordering,
which disproves the conjecture. \hfill$\square$

\paragraph{Scope and verification.}
The counterexample uses three effects and lies outside the source's
restricted class of sortable measurements: for instance,
$P_1-P_2=\operatorname{diag}(3/4,-3/4,0)$ is indefinite.
The conjecture explicitly concerns measurements beyond that class.
All effects and blocks are positive definite, so the failure does not
depend on zero effects or singularities.
The accompanying script checks the sequential product, all
normalizations, positivity and the squared-norm gap $3/64$ exactly.
All thirteen checks pass, including all six input orderings.

\paragraph{Task 2. Quantify the possible increase.}
Bound the increase of
$\max_i\|P_i-I_d/n\|_2$ under blockwise bistochastic dynamics, in terms
of the number of effects and the dimension.

\paragraph{Task 3. Find valid nonlinear monotones.}
Identify nonlinear quantities that are monotone on all quantum
measurements under the source's sequential product.

\paragraph{Task 4. Characterize the diagonal case.}
For simultaneously diagonal effects and blocks, determine when a
single ordering of the input effects controls every required partial
sum for a specified output ordering.

\paragraph{Task 5. Determine conditions restoring the conjecture.}
Find structural conditions beyond the source's sortable class that
ensure the asserted inequalities, and determine whether they are
preserved by blockwise bistochastic dynamics.

\begin{thebibliography}{9}
\bibitem{RZ}
A. Rico and K. \.{Z}yczkowski,
\emph{Discrete dynamics in the set of quantum measurements},
Journal of Physics A: Mathematical and Theoretical \textbf{57} (2024),
435302. Appendix C.1, Conjecture 1, equation (C.1);
\href{https://doi.org/10.1088/1751-8121/ad7dc2}{published article},
\href{https://arxiv.org/abs/2308.05835}{arXiv:2308.05835},
version 2, p.~28.
\end{thebibliography}
\end{document}
